\documentclass[11pt]{article}

\usepackage[a4paper,margin=1in]{geometry}
\usepackage{amsmath,amssymb,amsthm}
\usepackage{booktabs}
\usepackage{longtable}
\usepackage{enumitem}
\usepackage{hyperref}
\usepackage{tikz}
\usepackage{listings}
\usepackage{xcolor}
\lstset{
  basicstyle=\ttfamily\small,
  breaklines=true,
  frame=single,
  backgroundcolor=\color[gray]{0.97}
}

\title{ESS Burn-In Screening System: Modules A, B, and Safety-First Ensemble}
\author{}
\date{September 4, 2026}

\begin{document}
\maketitle

\begin{abstract}
This document describes a Python system for high-reliability electronic component screening in burn-in/ESS workflows.  
The project combines:
\begin{itemize}[nosep]
  \item Module A: dynamic and static anomaly detection on observed checkpoints;
  \item Module B: explainable 168-hour drift prediction from 0 h and 24 h measurements;
  \item Ensemble engine: deterministic safety-first fusion of Module A and Module B outputs at the 24-hour checkpoint.
\end{itemize}
The system is built as a set of reusable modules with CLI, API, and Streamlit dashboard interfaces, plus reproducible evaluation and synthetic benchmarks.
\end{abstract}

\tableofcontents
\newpage

\section{Project Scope and Objectives}

The project addresses latent defects that can evade static pass/fail limits but show anomalous drift behavior over time:
\begin{itemize}[nosep]
  \item Traditional static checks can miss failures that pass absolute bounds but deviate from peer behavior.
  \item Module A adds dynamic, peer-aware outlier detection.
  \item Module B predicts future drift from early checkpoints and flags unsafe forecasts.
  \item Ensemble combines both signals with conservative precedence for early reject decisions.
\end{itemize}

Primary objectives:
\begin{itemize}[nosep]
  \item Keep false negative rate (missed defects) extremely low.
  \item Provide traceable, explainable reasons for every classification.
  \item Allow safety-first integration in early-life screening workflows.
\end{itemize}

\section{Repository Organization}

\begin{longtable}{p{0.33\textwidth}p{0.62\textwidth}}
\toprule
\textbf{Path/Component} & \textbf{Purpose} \\
\midrule
\texttt{src/ess\_module\_a/} & Module A core engine, config, validation, API, CLI, detectors, statistics, feature engineering, evaluation, synthetic data generation. \\
\texttt{src/ess\_module\_b/} & Module B data adaptation, model training, regression engine, evaluation, API, CLI, feature construction. \\
\texttt{src/ess\_ensemble/} & Deterministic fusion engine, final evaluation logic, public exports. \\
\texttt{app.py} & Streamlit dashboard for interactive operation of Module A, Module B, and Ensemble sections. \\
\texttt{configs/parameters.yaml} & Module A engineering/static/robustness parameters, per-parameter limits and danger directions. \\
\texttt{configs/module\_b.yaml} & Module B model and safety configuration (e.g., ridge alphas, thresholds, risk multipliers). \\
\texttt{tests/} & Automated test suite for modules, adapters, APIs, evaluation, and edge cases. \\
\texttt{data/} & Input datasets and synthetic/acceptance data used for evaluation. \\
\texttt{artifacts/} & Generated model artifacts and precomputed acceptance reports. \\
\texttt{MODULE\_B.md} & Dedicated Module B technical documentation. \\
\texttt{README.md} & User-facing setup, API, and workflow documentation. \\
\texttt{Dockerfile}, \texttt{Dockerfile.module\_b}, \texttt{render.yaml} & Deployment support for UI and API services. \\
\bottomrule
\end{longtable}

\section{Data Contract}

\subsection{Long-form schema (primary)}
Required columns:
\begin{itemize}[nosep]
  \item \texttt{component\_id}, \texttt{lot\_id}, \texttt{part\_number}, \texttt{parameter}, \texttt{time\_h}, \texttt{value}, \texttt{unit}, \texttt{test\_condition\_id}
\end{itemize}

Required per-parameter context fields (from configuration):
\begin{itemize}[nosep]
  \item \texttt{temperature\_c}, \texttt{voltage\_v}, \texttt{test\_mode}
\end{itemize}

Optional engineering / label columns:
\begin{itemize}[nosep]
  \item \texttt{datasheet\_min}, \texttt{datasheet\_max}, \texttt{delta\_limit}, \texttt{actual\_value\_168h}, \texttt{is\_anomaly}, \texttt{defect\_type}
\end{itemize}

\subsection{Wide-organizer compatibility}
Adapters convert \texttt{Value\_0h}, \texttt{Value\_24h}, \texttt{Value\_168h} into the long form.  
Module B also accepts \texttt{Value\_168h} or \texttt{actual\_value\_168h} during training/evaluation.

\section{Module A: Dynamic Outlier Detection}

\subsection{High-level flow}
\begin{enumerate}
  \item Validate batch and convert units to canonical units.
  \item Build per-row history/historical features.
  \item Fit/load reference profile (lot-wise and global historical stats).
  \item Score each measurement using robust statistics and anomaly rules.
  \item Aggregate per-component worst-status and risk.
\end{enumerate}

\subsection{Validation and normalization}
Validation marks invalid records and logs structured issues:
\begin{itemize}[nosep]
  \item missing required columns/values,
  \item unsupported checkpoints,
  \item non-numeric values,
  \item unit or transform failures,
  \item duplicate records,
  \item missing required checkpoints up to \texttt{as\_of\_h}.
\end{itemize}
All values are converted by a dedicated unit converter into canonical units.

\subsection{Reference profile}
Reference fitting builds:
\begin{itemize}[nosep]
  \item historical distributions per \((\text{part},\text{parameter},\text{time},\text{condition})\),
  \item transformed-value counterparts,
  \item slope distributions by time transitions,
  \item lot-median baselines,
  \item optional unsupervised models:
  \begin{itemize}
  \item context-specific MinCovDet Mahalanobis models.
  \end{itemize}
\end{itemize}

\subsection{Robust statistical features and formulas}
For a value \(x\), reference median \(m\), MAD \(d\), and small constant \(\epsilon\):
\[
z_{\mathrm{robust}}=\frac{x-m}{\max(1.4826\,d,\epsilon)}
\]
This is applied on transformed values to both lot-local and historical groups.

Additional risk-aware quantities:
\[
\text{IQR distance}(x)=
\begin{cases}
\dfrac{x-Q_3}{IQR}, & x>Q_3\\
\dfrac{x-Q_1}{IQR}, & x<Q_1\\
0, & Q_1\le x\le Q_3
\end{cases}
\]
and direction-aware risk:
\[
R=
\begin{cases}
z, & \text{higher-is-bad}\\
-z, & \text{lower-is-bad}\\
|z|, & \text{two-sided}
\end{cases}
\]

\subsection{Rule categories}
Static checks produce immediate fail when spec bounds are broken:
\[
\text{margin} = \min(\text{spec\_max}-x,\;x-\text{spec\_min})
\]
Dynamic rules include robust-\(z\), IQR, and tail-percentile checks from lot and historical references.  
Each rule maps to warning/severe categories and reason codes.

\subsection{Status logic}
Let \(\text{StaticFail}\), \(\text{QualityIssues}\) be booleans; \(E=|\text{severe}|\), \(W=|\text{warning}|\), and evidence weight \(w=2E+W\).  
Then, for a valid row:
\[
\text{status}=\begin{cases}
\texttt{STATIC\_FAIL}, & \text{if static fails}\\
\texttt{RETEST\_REQUIRED}, & \text{if quality issues or missing references}\\
\texttt{QUARANTINE}, & \text{if severe extreme OR } w\ge 8 \text{ OR corroborated slope logic}\\
\texttt{MONITOR}, & w\ge 6\\
\texttt{NORMAL}, & \text{otherwise}
\end{cases}
\]
\[
\text{risk\_score}=\min\left(1,\max\left(\frac{\max(0,R_{\mathrm{lot}})}{z_{\mathrm{extreme}}},\frac{\max(0,R_{\mathrm{hist}})}{z_{\mathrm{extreme}}},\frac{p_{\mathrm{mah}}}{100}\right)\right)
\]
where \(p_{\mathrm{mah}}\) is the Mahalanobis percentile; static failure or major data issues force risk \(=1\).

\section{Module B: 168-hour Drift Predictor}

\subsection{Input restriction and leakage safety}
Module B intentionally uses only early-time data:
\[
\text{features are derived from } \{x_{0},x_{24}\}
\]
Future measurements (\(\ge 96\) h, \(\ge 168\) h targets) are ignored for forecasting input.

\subsection{Engineered features}
Feature vector:
\[
[x_{0},x_{24},\Delta_{0,24},s_{0,24},\text{relative\_drift},\log(x_{24}/x_0)]
\]
with slope \(\Delta_{0,24}/(24-0)\) and safe log handling with \(\epsilon\).

\subsection{Model candidates and selection}
For each context \((\text{part\_number},\text{parameter},\text{test\_condition\_id})\), candidates are:
\begin{itemize}[nosep]
  \item Ridge with \(\alpha\in\{0.01,0.1,1.0,10.0\}\),
  \item Huber regression.
\end{itemize}
All pipelines are standardized (StandardScaler + regressor).  
Lot-grouped \(K\)-fold CV (on lot IDs) provides out-of-fold MAE for each candidate:
\[
\mathrm{MAE}_{\text{candidate}}=\frac1N\sum_i|y_i-\hat y_i^{\mathrm{CV}}|
\]
Winner is min MAE, with deterministic tie preference for Ridge-style stability.

\subsection{Baselines}
Besides the model forecast \(\hat y_{168}^{\text{model}}\):
\[
\hat y_{168}^{\text{persist}}=x_{24}
\]
\[
\hat y_{168}^{\text{linear}} = x_{24} + 6(x_{24}-x_0)=7x_{24}-6x_0
\]

\subsection{Safety slope}
Direction-aware total drift rate:
\[
 s_{\mathrm{pred}} = \frac{\max(0,\mathrm{dir}(\hat y_{168}-x_0))}{168}
\]
where \(\mathrm{dir}(\cdot)\) applies danger direction (higher/lower/abs).

Safety-slope candidates:
\[
 s_{\delta} = \frac{|\Delta_{\mathrm{limit}}|}{168}
\]
\[
 s_{\mathrm{hist}}=\text{historical safe quantile}
\]
\[
 s_{\mathrm{guard}}=\max\left(0,\frac{\text{spec\_guard}-x_0}{168}\right)
\]
The binding safety slope is the minimum positive candidate (when available).

\subsection{Conservative interval and decision}
\[
 s_{\mathrm{conv}} = \frac{\max(0,\mathrm{dir}(x_{\text{conservative}}-x_0))}{168}
\]
with \(x_{\text{conservative}}\) adjusted by danger-side residual margin (e.g., + for higher-bad, - for lower-bad).  
Decision rules (priority):
\begin{itemize}[nosep]
  \item \texttt{STATIC\_FAIL} if any early measurement violates spec limits or if unsafe limit context exists.
  \item \texttt{RETEST\_REQUIRED} if limits/constraints missing/invalid.
  \item \texttt{EARLY\_REJECT} if
  \[
  s_{\mathrm{pred}}>s_{\mathrm{safe}}
  \]
  or linear sentinel confirms strongly unsafe acceleration.
  \item else \texttt{CONTINUE\_SCREENING}.
\end{itemize}
Risk is normalized against safety slope with conservative upper cap at 1.

\section{Ensemble: Safety-First Fusion}

\subsection{Precedence rules}
For each component:
\begin{enumerate}[label=\arabic*.]
  \item \textbf{REJECT\_EARLY} if Module A status is \texttt{QUARANTINE} or \texttt{STATIC\_FAIL} OR Module B decision is \texttt{EARLY\_REJECT} or \texttt{STATIC\_FAIL}.
  \item \textbf{RETEST\_REQUIRED} if any module output is missing or invalid.
  \item \textbf{MONITOR} if Module A is \texttt{MONITOR} and no hard reject.
  \item \textbf{CONTINUE\_SCREENING} otherwise.
\end{enumerate}

\subsection{Risk fusion}
\[
 r_E = 1-(1-r_A)(1-r_B)
\]
bounded in \([0,1]\). This is probability-style union risk that never decreases either module’s individual risk.

\section{Evaluation Design and Results}

\subsection{Metrics}
For binary screening outcome \(y\in\{\text{defect},\text{normal}\}\), decision \(d\in\{\text{action},\text{no action}\}\):
\[
\mathrm{Recall}=\frac{TP}{TP+FN},\qquad
\mathrm{Precision}=\frac{TP}{TP+FP},\qquad
\mathrm{FalseFlag}=\frac{FP}{FP+TN}
\]
Weighted screening cost:
\[
 C = 10\cdot FN + 1\cdot FP
\]
with false-negative multiplier from Module B config (10.0).

\subsection{Splitting}
Lot-safe evaluation splits lots into train / validation / test using fixed random seed and fraction policy so that no lot appears in multiple partitions.
\[
\text{train}\cup \text{validation}\cup \text{test}=\text{all lots},\qquad \text{disjoint partitions}
\]

\subsection{Observed benchmark snapshot (from checked-in artifacts)}
\paragraph{Module A acceptance report (snapshot)}
\begin{itemize}[nosep]
  \item Test: \(\text{defect recall}=1.0\), \(\text{false flag}=0.0204\).
  \item Validation: \(\text{defect recall}=1.0\), \(\text{false flag}=0.0340\).
\end{itemize}

\paragraph{Module B acceptance report (snapshot)}
\begin{itemize}[nosep]
  \item Test: \(TP=34, FN=1, FP=28, TN=537\),
  \[
  \mathrm{Recall}=0.9714,\quad \mathrm{Precision}=0.5484,\quad \mathrm{FalseFlag}=0.04956,\quad \mathrm{MAE}=0.4376
  \]
  \item Validation: \(TP=10, FN=0, FP=55, TN=543\), \(\mathrm{Recall}=1.0\), \(\mathrm{FalseFlag}=0.0797\), \(\mathrm{MAE}=0.2285\).
\end{itemize}

\paragraph{Ensemble acceptance report (snapshot)}
\begin{itemize}[nosep]
  \item Test: \(\text{component count}=600\), \(\text{RECALL}=0.9714\), \(FN=1\), \(FP=38\), \(\text{FalseFlag}=0.0673\), \(w\_cost=48\).
  \item Pattern counts: \texttt{BOTH\_MODULES\_REJECT}=11, \texttt{MODULE\_A\_ONLY\_REJECT}=5, \texttt{MODULE\_B\_ONLY\_REJECT}=51.
\end{itemize}

\section{Interfaces and Deployment}

\subsection{CLI}
\begin{itemize}[nosep]
  \item \texttt{ess-module-a}: generate/fit/score/evaluate/serve
  \item \texttt{ess-module-b}: fit/forecast/evaluate/serve
  \item \texttt{ess-ensemble}: lot-safe evaluate and emit final acceptance report
\end{itemize}

\subsection{FastAPI endpoints}
\begin{itemize}[nosep]
  \item Module A: \texttt{GET /v1/module-a/model-info}, \texttt{POST /v1/module-a/score-lot}, \texttt{GET /health}
  \item Module B: \texttt{GET /v1/module-b/model-info}, \texttt{POST /v1/module-b/forecast-lot}, \texttt{GET /health}
\end{itemize}

\subsection{Streamlit app}
Three interactive sections:
\begin{itemize}[nosep]
  \item Module A diagnostics and component evidence,
  \item Module B explanation and per-parameter forecast evidence,
  \item Ensemble final results, agreement matrix, fused risk, audit trail.
\end{itemize}

\section{Quality, Explainability, and Traceability}

Every output record preserves:
\begin{itemize}[nosep]
  \item per-parameter reason codes,
  \item risk score,
  \item decision labels,
  \item selected model and version metadata,
  \item forecast feature values and contributions (for linear/regression path),
  \item decision sources and validation issues.
\end{itemize}
Missing/invalid inputs are mapped to explicit retest requirements, not silent clears.

\section{Operational Notes}
\begin{itemize}[nosep]
  \item Synthetic metrics are prevalidated and useful for software acceptance; not equivalent to hardware qualification evidence.
  \item For production adoption, retrain artifacts on approved hardware lots and lock configuration versions.
  \item Any change to thresholds (e.g., recall bias, quantiles, guard band fraction) should be documented via artifact and config versioning.
\end{itemize}

\section{References to Config Files}
\begin{itemize}[nosep]
  \item \texttt{configs/parameters.yaml} (Module A defaults and thresholds),
  \item \texttt{configs/module\_b.yaml} (Module B model/safety hyperparameters),
  \item \texttt{artifacts/acceptance\_evaluation.json},
  \item \texttt{artifacts/module\_b\_acceptance\_evaluation.json},
  \item \texttt{artifacts/ensemble\_acceptance\_evaluation.json}.
\end{itemize}

\section{Conclusion}
This project integrates rule-based statistical screening, machine-learning-based drift forecasting, and conservative ensemble fusion to reduce escape risk with auditable behavior.  
The design intentionally favors traceability and safety over maximal model opacity, with explicit retest and reject precedence, lot-safe evaluation, and a full operational interface stack (CLI, API, UI).
\end{document}
